% II. Проектиране и програмна реализация. Описва реализираната структура и само
% решенията, които не са очевидни от подписите.
\section{Метод}
\label{sec:design}

\subsection{Структура и абстракция върху времето}

Системата е на четири слоя \figref{fig:layers}: изпълнителна среда, протоколно ядро, три
изпълними програми и слой за анализ върху маршрутните таблици и дневника със събития.

\begin{figure}[!htb]
\centering
\begin{tikzpicture}[node distance=0.9cm, scale=0.78, transform shape]
    \node[block, minimum width=3.4cm] (demo) {\code{waypoint-demo}\\сценарий и DOT};
    \node[block, minimum width=3.4cm, right=0.5cm of demo] (bench) {\code{waypoint-bench}\\измервания};
    \node[block, minimum width=3.4cm, right=0.5cm of bench] (live) {\code{waypoint-live}\\демон};

    \node[accent, minimum width=11.2cm, below=1.0cm of bench] (core)
        {Протоколно ядро: \code{Router}, \code{Lsdb}, \code{Graph}, \code{shortest\_paths}\\
         без обръщение към часовник и към сокет};

    \node[block, minimum width=11.2cm, below=0.8cm of core] (env)
        {Изпълнителна среда \code{Env}: \code{now}, \code{schedule}, \code{send}, \code{random}, \code{log}};

    \node[block, minimum width=5.4cm, below left=0.9cm and -5.6cm of env] (sim)
        {\code{Simulator} и \code{NodeEnv}\\виртуален часовник, опашка от\\събития, модел на връзките};
    \node[block, minimum width=5.4cm, below right=0.9cm and -5.6cm of env] (udp)
        {\code{UdpEnv} и \code{UdpSocket}\\\code{steady\_clock}, \code{select},\\Winsock или BSD сокети};

    \node[note, right=0.6cm of core, text width=3.0cm] (log)
        {Дневник със събития и\\анализ: сходимост,\\цикли, натоварване};

    \draw[line] (demo.south) -- ++(0,-0.35) -| (core.north);
    \draw[line] (bench) -- (core);
    \draw[line] (live.south) -- ++(0,-0.35) -| (core.north);
    \draw[line] (core) -- (env);
    \draw[line] (sim.north) -- ++(0,0.35) -| (env.south);
    \draw[line] (udp.north) -- ++(0,0.35) -| (env.south);
    \draw[line, dashed] (core.east) -- (log.west);
\end{tikzpicture}
\caption{Слоеве на системата. Стрелките сочат посоката на зависимостта.}
\label{fig:layers}
\end{figure}

Границата е структурна: заглавните файлове на платформата се включват единствено в
\code{src/udp.cpp}, така че нарушението ѝ се вижда като нов такъв ред в друг файл. Двата
режима изпълняват едни и същи обектни файлове на ядрото.

\subsection{Протоколно ядро}

Ядрото поддържа таблица със съседите, база с обявленията, индексирана по подател и с поредни
номера, и маршрутна таблица от дървото от най-кратки пътища върху графа, реконструиран от
базата. Графът е \term{насочен}, защото всеки маршрутизатор обявява само своята страна на
връзката, и преди изчислението се взема двупосочното му подмножество: ребро участва само ако и
двата му края са го обявили. Това правило е причината едностранният отказ да води до правилно
пренастройване, а не до постоянна черна дупка.

Разделението между база и маршрутна таблица е съществено за измерването: моментът, в който
възел научава за промяна, моментът, в който изчислява, и моментът, в който новата таблица
влиза в сила, са три различни момента. Инсталирането е отделна стъпка и може да бъде забавено
с параметър, защото интервалът между изчислението и влизането в сила е мястото, където
възникват преходните цикли.

\subsection{Симулатор, дневник и възпроизводимост}

Опашката е подредена по двойката (момент, пореден номер на вписване). Поредният номер е
съществен: без него две събития с еднакъв момент биха се изпълнявали в ред, зависещ от
контейнера. Всяка връзка се моделира като две независими посоки със закъснение, капацитет и
момент, до който посоката е заета, тоест моделът включва опашка и сериализация, не само
закъснение. Отказите са четири вида: отказ на връзка в двете посоки, отказ само в едната, отказ
на възел и възстановяване.

Всяко значимо събитие се записва като запис от шест полета: момент, възел, вид, отсрещен възел
и две числови полета. Анализът работи само върху този дневник и върху маршрутните таблици:
сходимостта се проверява чрез сравнение на таблицата на всеки възел с таблица, изчислена
независимо върху истинската топология, а преходните цикли чрез проследяване на следващите
стъпки за всяка двойка източник и цел. Двете проверки се извикват след всяко инсталиране, така
че моментите се засичат с точността на събитието. Възпроизводимостта се проверява, а не се
твърди: дневникът се свежда до 64-битова контролна сума и всеки сценарий се изпълнява два пъти.

\subsection{Формат по мрежата}

Всички цели числа се кодират явно, байт по байт, в мрежов ред (от най-значимия), по формата на
OSPFv2~\cite{rfc2328}: заглавие от 24 байта, нулева автентикация, контролна сума на пакета по
IP (без полето за автентикация) и контролна сума на Fletcher върху всяко обявление. Петте типа
пакети, които демонът използва (приветствие, описание на базата, заявка, обновяване и
потвърждение) съвпадат с петте типа на OSPFv2. Декодирането проверява обявената дължина и
двете контролни суми; обявеният брой обявления се проверява преди заделянето на памет. Златни
пакети, генерирани независимо и сверени за контролните суми, са част от тестовете.

\subsection{Методика на измерването}

Измерват се четири величини. \textbf{Време за сходимост}: интервалът от отказа до момента, в
който маршрутната таблица на всеки останал свързан възел съвпада точно с таблица, изчислена
независимо върху истинската топология, като съвпадението обхваща цената, следващите стъпки и
предшествениците, а не само достижимостта. \textbf{Преходни цикли}: интервали, в които
следващите стъпки на два или повече възела образуват затворен път. \textbf{Натоварване}:
съобщения и байтове след отказа. \textbf{Нестабилност}: брой цели, за които някой възел е
променил записа си. Времето в симулатора е виртуално, така че натовареността на машината
\textbf{не влияе на измерените стойности}. Резултатите са получени на Microsoft
Windows~11 Pro~N, версия~10.0.26200; AMD Ryzen~5~3600 6-Core Processor, 12 logical
processors; компилатор g++~15.2.0 MinGW, CMake~4.3.2 и Ninja~1.13.2 в режим
\code{Release}. Командата от корена на хранилището е
\code{.\textbackslash{}build\textbackslash{}waypoint-bench -{}-csv results.csv}. Времената в
таблиците са милисекунди виртуално време, не стенен часовник.

Топологиите са четири семейства: пръстен, който дава най-голям диаметър и е песимистичният
случай; решетка със средна свързаност; пълна мрежа с диаметър единица и максимално натоварване;
и дебело дърво, в което пътищата с равна цена са правило. За зависимостта от размера са взети
пръстени с 4 до 32 възела, тоест диаметър от 2 до 16; за таймерите пръстен с 12 възела; за
плътността четири топологии с по 12 възела и средна степен от 2,00 до 11,00; а за вида на
отказа пръстен, решетка четири на четири и дебело дърво с $k$ равно на четири.

Стойностите по подразбиране са: поддържане 1 s, отпадане 4 s (отношението по RFC
2328~\cite{rfc2328}), повторно изпращане 1 s, изчакване преди изчисление 100 ms, закъснение
при инсталиране 0, закъснение по връзка 1 ms и капацитет 100 Mbit/s. В експеримента за
таймерите поддържането приема от 250 ms до 4 s, а отпадането е четирикратно. В експеримента за
циклите всеки маршрутизатор получава закъснение при инсталиране от 0 до 200 ms, изтеглено от
зърното, защото записването на нова таблица отнема време. \textbf{Ако това закъснение е нула за
всички, преходен цикъл е невъзможен по построение}, и измереният нула би бил свойство на модела.

Всяка комбинация се изпълнява по десет пъти със зърна от 1000 до 1009, така че фазата между
последното приветствие и отказа, единственият източник на разсейване при виртуално време, да
бъде извадена достатъчно на брой пъти. Всяко изпълнение има загряване от 40 s, отказ на 41-ата
секунда и прозорец от 60 s. Резултат е валиден само ако мрежата е била сходима преди отказа,
ако е сходима отново в рамките на прозореца и ако повторението със същото зърно дава
\textbf{идентична контролна сума на дневника}.
